\paragraph{Cash classifications can overlap without revealing their members.}
The source notes make the overlap concrete. In the retained 2024 annual filing,
Note~7 on PDF page~42 (printed page~39) divides the USD~15.8 billion State Aid
escrow into USD~2.6 billion classified as cash and cash equivalents and
USD~13.2 billion classified as current marketable securities. Note~4 on PDF
page~40 (printed page~37) also states that the restricted cash is included in
cash and cash equivalents. Restricted status therefore cannot define a
disjoint addition to that balance-sheet caption. The corresponding tagged
cash disclosure retains its State Aid contingency dimension; it is not an
unqualified total of every restriction.

To express what the two cash totals can tell us, let $B$ denote the population
under the balance-sheet caption and $F$ the population under the cash-flow
caption, at one date. Assume for this comparison that both use the same
nonnegative additive monetary valuation. Write $b$ for the value belonging
only to $B$, $f$ for the value belonging only to $F$, and $i$ for the value in
their intersection. These are aggregate values of membership classes, not
identified bank accounts. If $C_B$ and $C_F$ are the two measured totals, the
disjoint decomposition gives
\[
 C_B=b+i,\qquad C_F=f+i,\qquad C_F-C_B=f-b,
 \qquad b,f,i\geq0.
\]
The common value cancels on subtraction. Treating the 2023 displayed amounts
as exact solely for this illustration, the difference of USD~772 million
therefore fixes $f-b$. Both $(b,f,i)=(0,772,29965)$ and
$(100,872,29865)$ satisfy the disclosed totals. An additional assumption
$B\subseteq F$ forces $b=0$ and hence $f=772$ under that valuation, but the
subtraction itself supplies no such membership evidence and does not identify
the financial nature of $f$.

Equal totals do not settle the matter either. At the 2024 displayed amount
of USD~29,943 million, both $(0,0,29943)$ and $(100,100,29843)$ satisfy the
two equations. Moreover, even zero exclusive value does not prove that the
exclusive set is empty when zero-value members are possible. The monetary
constraints preserve this distinction between equal values and equal members.
The 2023 annual filing independently displays the same earlier cash totals;
that cross-edition agreement does not identify the missing members. No exact
classification of the USD~772 million difference was located in the bounded
inspection of the retained annual statements, cash notes and relevant tagged
calculation arcs.

The optional \texttt{ConditionalPopulationOverlap/1} lowers two through four
declared populations to the existing joint-measurement engine. It retains
every nonempty Boolean membership class, the original source dimensions,
explicit bounds and proposed projection premises. It requires a common
additive valuation; signed or differently measured quantities need a separate
mapping. Bounds derived from displayed precision also remain proposed. For
example, a lower-closed, upper-open nearest-rounding interpretation places
the USD~2.6 billion cash disclosure in $[2550,2650)$ million; it does not turn
the displayed number into an exact underlying balance.

Under that rounded example, requiring the disclosed escrow cash both to belong
to $B$ and to be disjoint from $B$ is infeasible. This rejects the proposed
classification, not the filing. Eleven focused tests, 48 existing regression
tests, ten complete exchanges and four separately checked rational witnesses
cover the feasible alternatives, incompatible classifications, missing premises,
context mismatches, resource limits and forged results. The implementation
preserves unresolved membership and supplies no account postings or executable
company opening. The evidence and exact source ledger are retained at
\begin{quote}\small\ttfamily
data/accounting-master/\\
cash-population-reconciliation-20261025/\\
population-overlap-20261025/
\end{quote}
